\section{结论与讨论}

在本工作中，我们给出了一个具体的处方，用于非微扰地重正化准 PDF 计算所需的矩阵元。
所采用的方案是 RI$'$，随后转换到 $\MSb$ 方案并演化到 2 GeV；这一步以单圈微扰完成。
我们已经论证，该重正化条件恰当地处理了矩阵元中存在的两类发散：标准的对数发散，
以及含 Wilson 线的非局域算符特有的幂发散。此外，我们给出了消除混合的重正化条件，
适用于与非极化准 PDF——它与 twist-3 标量算符相混合——的情形。

\medskip
我们还演示了所提出处方在螺旋度准 PDF 上的实施，并给出了相应的重正化矩阵元。
这使我们得以得出如何使计算更加稳健的结论，而这正是本工作的主要成果。
\begin{itemize}
\item 首先，一项具有受控系统不确定性的计算的一个关键要素，是在格点微扰论框架内扣除单圈格点伪影。
循着文献~\cite{Alexandrou:2015sea} 的思路，我们目前正在计算 $\mathcal{O}(g^2\,a^\infty)$ 伪影，
并将把它们从 $Z$ 因子的非微扰估计值中扣除。这样，重正化函数（特别是对``平行''动量）中出现的
大截止效应将在很大程度上得到避免。
\item 其次，从 RI$'$ 重正化方案到 $\MSb$ 方案的转换因子很可能带有可观的高阶修正，它们（除其他作用外）
是重正化矩阵元实部在大 Wilson 线长度处变为负值这一非物理特征的成因。该转换因子的两圈计算
预期可以在足够程度上解决这一问题。我们所做的数值实验表明，转换因子由两圈贡献带来的自然改变
（即大约 10-20\%，近似为所考虑标度处的 $\alpha_s$）应当足以压制这一不想要的效果。
转换因子到两圈的微扰计算相当繁难，将另行发表。
\end{itemize}

\medskip
总而言之，本工作提出的重正化程序连同正在推进的未来改进，将为 Wilson 线费米子算符的
重正化函数提供可靠的估计。如此得到的重正化矩阵元可以用作输入来计算准 PDF 并将其匹配到光前 PDF，
这是整个方法的主要目标。除了本工作讨论的螺旋度情形之外，我们将以类似的方式处理横向性 PDF。
对于非极化情形，则需要考虑与标量算符的混合，正如这里所解释并数值演示的那样。
通过这项工作，我们提出并讨论了准 PDF 的完整重正化程序——在此之前它一直是一个重大的不确定性来源——
并构成了连接格点 QCD 结果与光锥 PDF 的关键里程碑。
